\chapter{Introduction}

As the digital energy field is expanding quickly, The amount of data stored is growing exponentially. With this surge, the role of data analytics has become a hugely important through the whole subsurface energy fields sectors (i.e. upstream and downstream). As the fields of big data analytics, artificial intelligence and machine learning grow, companies need data analysis in various forms and sectors through their chain. This makes potential use of data mining methodologies to learn from experts how they predict well failures. The performance can be dramatically improved by successful failure predictions either by adjusting operating parameters to forestall failures or by scheduling maintenance to reduce unplanned repairs and to minimize downtime. In the smart subsurface energy fields, the measurements from fields are used for efficient and integrated management. Hence machine learning and data analytics researches increased drastically.

Machine learning in the oil and gas industry can be divided into three main categories. Those are either extracting knowledge from data in real-time form, predictive analytics and production optimization through real data. As aforementioned, the oil and gas industry generates enormous data. However, there is no full usage of that data. With the help of machine learning, this data can be processed to make real time diagnostic applications. Furthermore, machine learning can help to make predictions regarding how oil production from wells. Also, it could help in optimizing production through identifying the most efficient ways to set up production systems. Such applications can include advisory systems for injection rates, pumping stroke per minutes and anything determining optimal well spacing, deciding where to drill, and testing different techniques for fracking.

These three categories exhibit how machine learning can help oil and gas companies make better decisions today and plan for the future based on past data. This will be critical as they face the challenges of improving their workflow while adapting to new technologies. From energy systems perspective, there are several areas where machine learning can be applied to help improve oil and gas industry workflows, including: Real-time drilling, reservoir engineering, Oil and gas production and procurement, downtime prevention and well testing. Etc

The above application areas include a wide range of tasks that involve oil and gas industry employees. Machine learning can help these employees become more efficient at their jobs, help them make better decisions, and improve the quality of products they produce. However, the transition will not happen overnight. As the oil and gas industry prepares to incorporate machine learning, it is important to recognize the challenges involved. The first significant challenge is the oil and gas industry’s need to hire data scientists qualified to extract knowledge from industry data. This may be difficult for companies that already have a shortage of workers or companies in regions where there aren’t many people with data science skills. 

Another challenge is the large amount of data oil and gas companies deal with daily. This means that the machines doing the analysis will need to be very powerful to process it all in a reasonable amount of time. Finally, training the machine learning algorithms to recognize valuable data patterns for oil and gas companies will require a great deal of time and effort. However, the payoff of accurate models to make data-backed decisions will be worth it for oil and gas companies.

The realization that machine learning techniques can be successfully deployed for oil and gas companies is still occurring across the industry. As more energy companies investigate ML use cases, we will see an increase in its adoption. This will certainly change how oil and gas companies operate, and it is just one example of how data science is changing how we do business.

Oil and gas companies are already using machine learning to help with prediction models, produce better results from their resources, and optimize their processes. Machine learning has the power to revolutionize the oil and gas industry. This forward momentum is crucial for the sector solely responsible for keeping the rest of society moving.

In the coming years, we will see more and more ML technologies used by oil and gas companies around the world. In fact, the future of the oil and gas industry can be found in the convergence of many different technologies, including machine learning. Companies who adapt early can improve their resource optimization and gain a more detailed look at their progress.


This research is prepared to develop applications on the whole framework of data analytics applications on the subsurface energy systems industry. Applications on three stages of data analysis are developed. These three stages are descriptive, predictive and prescriptive analytics. Firstly a virtual flow meters on the electrical submersible pump as an application on descriptive bases. Secondly, predictive maintenance of the electrical submersible pump is developed as an application on predictive analytics. On prescriptive analytics basis, the reinforcement learning is used to provide an optimum policy of steam injection rate.

This thesis is organized into eight chapters. Chapter 1 “Introduction” is prepared to give a general overview about the subject of the present study. Chapter 2 “Literature Review” presents a review for the data analysis and relevant literature on oilfield applications of real time, predictive and prescriptive analytics on petroleum production  systems. It also discussed the applicability of these various approaches. Then, concluding remarks. Chapter 3 “Statement of the Problem, Thesis Objectives and Methodology” describes the statement of the problem and the main objectives of the study and workflow pipeline for each application. Chapter 4 “virtual flow metering” Discussing the data driven multiphase flow estimation through artificial lift systems with focus on electrical submersible pump. Chapter 5 “Predictive modeling” discuss predive maintenance of the electrical submersible pumps by predicting pre-workover events using principle component analysis pipelined with extreme gradient boosting. Chapter 6 “prescriptive analytics” discuss the use of actor critic reinforcement learning to lead decision-making in energy systems optimization and introduces also a proof of value on steam injection optimization. Chapter 7 “Results and Discussion” shows the proposed models accuracy and various matrices used for testing to validate models. Chapter 8 “Conclusion and Recommendations” concludes the thesis results and findings. In addition, it presents the recommendations for the future research work in this topic.
